% Figure: the deviatoric (pi) plane in principal-stress space. The
% hydrostatic axis sigma1 = sigma2 = sigma3 runs out of the plane; the
% three principal axes project onto the plane at 120 degrees. The von
% Mises criterion is a circle on this plane (yield depends only on the
% deviatoric stress magnitude); the Tresca criterion is the inscribed
% regular hexagon. A purely hydrostatic stress sits at the origin and
% never yields. The figure shows why von Mises is rotationally symmetric
% on the deviatoric plane while Tresca has corners.
\begin{figure}[htbp]
\centering
\begin{tikzpicture}[scale=1.0]
% radius of the von Mises circle on this plane
\def\rad{2.2}
% von Mises circle
\draw[engblue, very thick] (0,0) circle[radius=\rad];
% the three projected principal axes at 90, 210, 330 degrees
\foreach \ang/\lab in {90/{$\sigma_1$}, 210/{$\sigma_2$}, 330/{$\sigma_3$}}{
  \draw[enggray, thick, -{Latex}] (0,0) -- (\ang:{\rad+0.7});
  \node[enggray, font=\small] at (\ang:{\rad+1.05}) {\lab};
}
% Tresca hexagon: regular hexagon inscribed in the circle, with vertices
% on the projected principal axes (vertices at 90, 150, 210, 270, 330, 30).
\draw[engred, very thick]
  (90:\rad) -- (150:\rad) -- (210:\rad) -- (270:\rad) -- (330:\rad) -- (30:\rad) -- cycle;
% origin = hydrostatic point
\filldraw[engblack] (0,0) circle[radius=1.6pt];
\node[engblack, font=\scriptsize, anchor=north east] at (-0.05,-0.05)
  {hydrostatic point};
% labels for the two loci
\node[engblue, font=\small, anchor=south west] at (30:\rad) {};
\node[engblue, font=\small] at (60:{\rad+0.0}) {};
\draw[engblue, -{Latex}] (1.0,2.6) -- (0.78,{0.78*\rad/1.1});
\node[engblue, font=\scriptsize, anchor=south] at (1.0,2.6) {von Mises (circle)};
\draw[engred, -{Latex}] (-2.7,-1.7) -- (-1.55,-1.55);
\node[engred, font=\scriptsize, anchor=east] at (-2.7,-1.7) {Tresca (hexagon)};
% a radial line marking the deviatoric stress magnitude
\draw[engorange, very thick, -{Latex}] (0,0) -- (18:{\rad});
\node[engorange, font=\scriptsize, anchor=west] at (18:{0.55*\rad}) {$\sqrt{2/3}\,\sigma_{\text{vM}}$};
\end{tikzpicture}
\caption[The deviatoric plane and the two yield loci]{The von Mises and
Tresca criteria viewed on the deviatoric (octahedral) plane in
principal-stress space, the plane perpendicular to the hydrostatic axis
$\sigma_1 = \sigma_2 = \sigma_3$. The three principal axes project onto
the plane at $120^{\circ}$ to one another. Von Mises is a circle: yield
depends only on the magnitude of the deviatoric stress, the radial
distance from the hydrostatic point at the origin, and not at all on
direction, so the criterion is rotationally symmetric here. Tresca is
the regular hexagon inscribed in the circle, touching it at the six
points of uniaxial tension and compression and falling inside it
elsewhere, which is the geometric statement that Tresca is the more
conservative criterion. A purely hydrostatic stress maps to the origin
and lies inside every locus, the picture's version of the fact that
hydrostatic pressure alone never causes ductile yield.}
\label{fig:vol03ch03:von-mises-octahedral}
\end{figure}
